%%%%%%%%%%%%%%%%%
% This is an sample CV template created using altacv.cls
% (v1.7.2, 28 Aug 2024) written by LianTze Lim (liantze@gmail.com), based on the
% CV created by BusinessInsider at http://www.businessinsider.my/a-sample-resume-for-marissa-mayer-2016-7/?r=US&IR=T
%
%% It may be distributed and/or modified under the
%% conditions of the LaTeX Project Public License, either version 1.3
%% of this license or (at your option) any later version.
%% The latest version of this license is in
%%    http://www.latex-project.org/lppl.txt
%% and version 1.3 or later is part of all distributions of LaTeX
%% version 2003/12/01 or later.
%%%%%%%%%%%%%%%%

%% Use the "normalphoto" option if you want a normal photo instead of cropped to a circle
% \documentclass[10pt,a4paper,normalphoto]{altacv}

\documentclass[10pt,a4paper,ragged2e,withhyper]{altacv}
%% AltaCV uses the fontawesome5 and simpleicons packages.
%% See http://texdoc.net/pkg/fontawesome5 and http://texdoc.net/pkg/simpleicons for full list of symbols.

% Change the page layout if you need to
\geometry{left=1.25cm,right=1.25cm,top=1.5cm,bottom=1.5cm,columnsep=1.2cm}

% The paracol package lets you typeset columns of text in parallel
\usepackage{paracol}


% Change the font if you want to, depending on whether
% you're using pdflatex or xelatex/lualatex
% WHEN COMPILING WITH XELATEX PLEASE USE
% xelatex -shell-escape -output-driver="xdvipdfmx -z 0" mmayer.tex
\iftutex
  % If using xelatex or lualatex:
  \setmainfont{Lato}
\else
  % If using pdflatex:
  \usepackage[default]{lato}
\fi

% Change the colours if you want to
\definecolor{VividPurple}{HTML}{3E0097}
\definecolor{SlateGrey}{HTML}{2E2E2E}
\definecolor{LightGrey}{HTML}{666666}
% \colorlet{name}{black}
% \colorlet{tagline}{PastelRed}
\colorlet{heading}{VividPurple}
\colorlet{headingrule}{VividPurple}
% \colorlet{subheading}{PastelRed}
\colorlet{accent}{VividPurple}
\colorlet{emphasis}{SlateGrey}
\colorlet{body}{LightGrey}

% Change some fonts, if necessary
% \renewcommand{\namefont}{\Huge\rmfamily\bfseries}
% \renewcommand{\personalinfofont}{\footnotesize}
% \renewcommand{\cvsectionfont}{\LARGE\rmfamily\bfseries}
% \renewcommand{\cvsubsectionfont}{\large\bfseries}

% Change the bullets for itemize and rating marker
% for \cvskill if you want to
\renewcommand{\cvItemMarker}{{\small\textbullet}}
\renewcommand{\cvRatingMarker}{\faCircle}
% ...and the markers for the date/location for \cvevent
% \renewcommand{\cvDateMarker}{\faCalendar*[regular]}
% \renewcommand{\cvLocationMarker}{\faMapMarker*}


% If your CV/résumé is in a language other than English,
% then you probably want to change these so that when you
% copy-paste from the PDF or run pdftotext, the location
% and date marker icons for \cvevent will paste as correct
% translations. For example Spanish:
% \renewcommand{\locationname}{Ubicación}
% \renewcommand{\datename}{Fecha}


%% Use (and optionally edit if necessary) this .tex if you
%% want to use an author-year reference style like APA(6)
%% for your publication list
% \input{pubs-authoryear.tex}

%% Use (and optionally edit if necessary) this .tex if you
%% want an originally numerical reference style like IEEE
%% for your publication list
\input{pubs-num.tex}

%% sample.bib contains your publications
\addbibresource{sample.bib}

\begin{document}
\name{Jescelin Hartin}
\tagline{Passionate Teacher}
% Cropped to square from https://en.wikipedia.org/wiki/Marissa_Mayer#/media/File:Marissa_Mayer_May_2014_(cropped).jpg, CC-BY 2.0
%% You can add multiple photos on the left or right
%\photoR{2.5cm}{mmayer-wikipedia-cc-by-2_0}
% \photoL{2cm}{Yacht_High,Suitcase_High}
\personalinfo{%
  % Not all of these are required!
  % You can add your own with \printinfo{symbol}{detail}
  \email{jescelinhartin@gmail.com}
  \phone{+971-501886350}
%  \mailaddress{Sharjah, UAE}
  \location{Sharjah, UAE}
 % \homepage{marissamayr.tumblr.com}
  % \twitter{@marissamayer}
 % \xtwitter{@marissamayer}
  \linkedin{jescelin-hartin   }
%   \github{github.com/mmayer} % I'm just making this up though.
%   \orcid{0000-0000-0000-0000} % Obviously making this up too.
  %% You can add your own arbitrary detail with
  %% \printinfo{symbol}{detail}[optional hyperlink prefix]
  % \printinfo{\faPaw}{Hey ho!}
  %% Or you can declare your own field with
  %% \NewInfoFiled{fieldname}{symbol}[optional hyperlink prefix] and use it:
  % \NewInfoField{gitlab}{\faGitlab}[https://gitlab.com/]
  % \gitlab{your_id}
	%%
  %% For services and platforms like Mastodon where there isn't a
  %% straightforward relation between the user ID/nickname and the hyperlink,
  %% you can use \printinfo directly e.g.
  % \printinfo{\faMastodon}{@username@instace}[https://instance.url/@username]
  %% But if you absolutely want to create new dedicated info fields for
  %% such platforms, then use \NewInfoField* with a star:
  % \NewInfoField*{mastodon}{\faMastodon}
  %% then you can use \mastodon, with TWO arguments where the 2nd argument is
  %% the full hyperlink.
  % \mastodon{@username@instance}{https://instance.url/@username}
}

\makecvheader

%% Depending on your tastes, you may want to make fonts of itemize environments slightly smaller
\AtBeginEnvironment{itemize}{\small}

%% Set the left/right column width ratio to 6:4.
\columnratio{0.6}

% Start a 2-column paracol. Both the left and right columns will automatically
% break across pages if things get too long.
\begin{paracol}{2}
  \cvsection{Summary}
  \begin{itemize}
   \item Dedicated and passionate Teacher,  fostering a positive and engaging learning environment. Skilled in assisting with lesson planning, classroom management, and individualized student support to enhance learning outcomes. Adept at creating interactive activities, maintaining a structured classroom, and providing one-on-one attention to students needing additional guidance.  \end{itemize}


\divider

\cvsection{Experience}

\cvevent{Teacher}{Everwin Vidhyashram}{Apr 2019 -- Dec 2019}{Chennai, India}
\begin{itemize}
\item Planning and implementing age-appropriate lessons aligned with the curriculum.
\item Foster a positive and inclusive classroom environment by promoting student engagement and participation.
\item Organize and prepare instructional materials, classroom activities, and visual learning aids.
\item Communicate with parents and guardians to provide updates on student progress and classroom activities.
  \end{itemize}

\divider

\cvevent{Internship Trainee}{CSI.St.Pauls Higher Secondary School}{Aug 2018 -- Dec 2018}{Chennai, India}
\begin{itemize}
\item Assist the lead teacher in planning age-appropriate lessons aligned with the curriculum.
\item Support students individually and in small groups to reinforce learning concepts and encourage academic growth.
\item Help manage classroom behavior using positive reinforcement strategies to ensure a structured and supportive learning space.
\item Assist with grading, attendance tracking, and maintaining student records.
\end{itemize}


\cvevent{Project Intern}{Siddha Central Research Insitute}{Nov 2016 -- Feb 2017}{Chennai, India}
\begin{itemize}
\item Worked on a Internship project titled “HPTLC Study \& In-vitro Antioxidant Assay of \‘Parutti Vitai Curanam\’ \– A Siddha Drug” under the guidance of Dr. R. Ganesan, Assistant Director (Biochemistry), SCRI, Chennai.
\item Developed the HPTLC profile of the hydroalcoholic extract of Parutti Vitai Curanam.
\item Analyzed and determined the polyphenol content, flavonoid concentration, and pro-oxidant effect of the extract.
\end{itemize}

% \divider

% \cvevent{Product Engineer}{Google}{23 June 1999 -- 2001}{Palo Alto, CA}

% \begin{itemize}
% \item Joined the company as employe \#20 and female employee \#1
% \item Developed targeted advertisement in order to use user's search queries and show them related ads
% \end{itemize}



% use ONLY \newpage if you want to force a page break for
% ONLY the currentc column





%% Switch to the right column. This will now automatically move to the second
%% page if the content is too long.
\switchcolumn

\cvsection{Education}
\cvevent{Bachelor of Education}{CSI.Bishop Newbigin College of Education}{June 2017 -- June 2019}{}
\cvevent{M.Sc Bio-Chemistry}{Valliammal College  For Women }{May 2015 -- May 2017}{}
\cvevent{B.Sc Bio-Chemistry }{Valliammal College For Women }{May 2012 -- May 2015}{}


\cvsection{Awards \& Honors}

\cvachievement{\faTrophy}{Bishop Sundar Clarke Gold Medal}{For Best Academic Record during 2017-2019}

\divider

\cvachievement{\faTrophy}{Bishop Lesslie Newbigin Gold Medal}{For securing highest marks in psychology of Learners and Learning}

\divider

\cvachievement{\faTrophy}{University Rank Holder}{Issued by University of Madras}

\divider


\cvsection{Strengths}
% Don't overuse these \cvtag boxes — they're just eye-candies and not essential. If something doesn't fit on a single line, it probably works better as part of an itemized list (probably inlined itemized list), or just as a comma-separated list of strengths.

% The `ragged2e` document class option might cause automatic linebreaks between \cvtag to fail.
% Either remove the ragged2e option; or 
% add \LaTeXraggedright in the paragraph for these \cvtag
{\LaTeXraggedright
\cvtag{Classroom Behavior Management}
\cvtag{Interactive Teaching Methods}
\cvtag{Encouraging Critical Thinking \& Creativity}
\cvtag{Multicultural \& Inclusive Teaching}
\cvtag{Classroom Organization \& Resource Management}
\par}



%\divider
\cvsection{Publications}

%% Specify your last name(s) and first name(s) as given in the .bib to automatically bold your own name in the publications list.
%% One caveat: You need to write \bibnamedelima where there's a space in your name for this to work properly; or write \bibnamedelimi if you use initials in the .bib
%% You can specify multiple names, especially if you have changed your name or if you need to highlight multiple authors.
\mynames{Jescelin Hartin}
%% MAKE SURE THERE IS NO SPACE AFTER THE FINAL NAME IN YOUR \mynames LIST

\nocite{*}

\printbibliography[heading=pubtype,title={\printinfo{\faFile*[regular]}{Journal Articles}}, type=article]

\divider

\end{paracol}

\end{document}